\section{When every height scale diverges}
\label{sec:high}

We now prove the auxiliary sequence theorem when all its scales tend to
infinity.  Proposition~\ref{prop:tangent} supplies positivity on every
subproduct through the marked points.  To use this positivity in a height
inequality, we first bound the intersection numbers from below in terms
of the degrees of those subproducts.  The resulting bound will be uniform
over a fixed neighborhood of the real projector.  This uniformity lets us
choose the rational approximation only after fixing the multiplicative
constant in the height inequality.

\subsection{A uniform lower bound for intersections}

Fix symmetric very ample line bundles $L_\ell$ on the abelian varieties
$A_\ell$, and a symmetric very ample line bundle $L$ on $\cA$.  We may
replace the $L_\ell$ by powers so that their embeddings are projectively
normal.  Give $L$ its invariant Chern form $\omega_L$.
For an irreducible subproduct $Z=\prod_\ell Z_\ell$, write
\[
 d_\ell=\dim Z_\ell,\qquad d=\sum_\ell d_\ell,
 \qquad D_Z=\prod_\ell\deg_{L_\ell}Z_\ell.
\]
We use the function $P_Z(g)=\int_Z(g^*\omega_L)^d$ defined in the
preceding section.  A zero-dimensional factor is a point over $\Qbar$;
it contributes degree one and dimension zero.

\begin{lemma}\label{lemma:intersection-lower-bound}
Let $\mathcal B$ be a fixed closed box with rational endpoints in the
coordinates of an integral basis of $\End(\cA)$.  There are positive
integers $\theta,\omega_1$, depending only on the fixed ambients,
polarizations, and box, with the following property.  For every
irreducible subproduct $Z\subseteq\cA$ such that $P_Z$ is strictly
positive everywhere on $\mathcal B$,
\begin{equation}\label{eq:intersection-lower-bound}
 P_Z(g)\ge \theta^{-1}D_Z^{-\omega_1}
 \qquad (g\in\mathcal B).
\end{equation}
\end{lemma}

\begin{proof}
Put $n=\dim\cA$ and $q=\rank_{\mathbb Z}\End(\cA)$, and choose
an integral basis $e_1,\ldots,e_q$.  The polynomial
\[
 p_Z(x_1,\ldots,x_q)
   =P_Z\left(\sum_{j=1}^q x_je_j\right)
\]
has degree $2d\le 2n$: pullback of a Hermitian form is quadratic in
the real coordinates of the map.  At integral arguments it is an
integer, because it is the top intersection number of the pullback of
$L$ by an actual endomorphism.

Set $K=2n$.  In the multivariate Newton expansion,
\[
 p_Z(x)=\sum_{|\alpha|\le K}(\Delta^\alpha p_Z)(0)
                 \prod_{j=1}^q\binom{x_j}{\alpha_j},
\]
all finite differences are integers.  Since
$\prod_j\alpha_j!$ divides $K!$, every coefficient of $p_Z$ has
denominator dividing $Q=K!$.  This denominator is independent of $Z$.

The coefficients also have absolute value at most $C D_Z$ for a fixed
$C$.  Indeed, on the fixed interpolation grid $\{0,\ldots,K\}^q$,
the semipositive forms $g^*\omega_L$ are bounded above by a fixed
multiple of
$\Omega=\sum_\ell\operatorname{pr}_\ell^*\omega_{L_\ell}$.
Consequently
\[
 0\le p_Z(x)\le C_0^d\int_Z\Omega^d
   =C_0^d\frac{d!}{\prod_\ell d_\ell!}D_Z.
\]
Fixed-grid interpolation gives the coefficient bound.  There are only
finitely many possible dimension tuples, so the constants are uniform.
These integrations can be taken over the regular locus, and the same
argument therefore includes singular subvarieties.

Here is the real-algebraic step which turns these bounds into a lower
bound for a positive minimum.  Let $p_c(x)$ be the universal polynomial
of degree at most $K$, whose coefficients form the free parameter vector
$c$.  Apply quantifier elimination over the real numbers to
\[
 \exists x\,\bigl(x\in\mathcal B\ \text{and}\ t=p_c(x)\bigr),
\]
leaving $c$ and $t$ free.  The resulting formula is a Boolean combination
of sign conditions on finitely many polynomials
$R_j(c,t)\in\mathbb Z[c,t]$; denominators from the rational box have
been cleared.  Quantifier elimination with free variables, including
its coefficient-ring formulation, is recalled in
\cite[arXiv version, Section~2.5.2, Theorem~2.27]{Basu2017}.

Specialize to the coefficient vector $c(Z)$, and put
$m_Z=\min_{\mathcal B}p_Z>0$.  At least one specialized polynomial
$R_j(c(Z),t)$ which is not identically zero must vanish at $m_Z$.
Otherwise the signs of every nonzero specialized polynomial would be
constant in a neighborhood of $m_Z$, and the identically zero ones
would have constant sign zero there.  The formula describing the range
would then be locally constant.  This is impossible: it holds at
$m_Z$ and fails immediately to its left.

Let $b=\max_j\deg_c R_j$ and $h=\max_j\deg_t R_j$.
Multiplication by $Q^b$ makes each specialization an integer polynomial,
of degree at most $h$ and coefficient size at most $C_1D_Z^b$.
Choose a nonzero one which vanishes at $m_Z$ and remove its maximal
power of $t$.  Its constant coefficient is now a nonzero integer.  If
$m_Z\le1$, its root equation implies
\[
 1\le \sum_{k=1}^h |a_k|m_Z^k
     \le hC_1D_Z^b m_Z.
\]
For $m_Z>1$ a lower bound is immediate.  Enlarging fixed integer
constants gives \eqref{eq:intersection-lower-bound}.  The case $d=0$,
where $p_Z=1$, is included directly.
\end{proof}

\subsection{The height inequality and its constants}

The height inequality follows the line of Vojta's curve method
\cite{Vojta1991}, Faltings' abelian-variety method \cite{Faltings1991},
and R\'emond's generalized projective inequality
\cite{Remond2005Vojta}. We use Dill's version with separate height
thresholds for the factors; his introduction credits Ange with the
preceding extension to different projective varieties
\cite[Introduction]{Dill2020}. We record its precise specialization
below; the independence of its multiplicative constant from coefficient
heights is essential to our approximation argument.  All point heights in this subsection are
absolute logarithmic projective heights with the maximum norm at every
place.  The notation $h(X)$ in the following statement denotes the
projective variety height with the normalization in the setup of
\cite[Theorem~1.1]{Dill2020}.

\begin{theorem}[Dill]\label{thm:dill-form}
Let $X_1,\ldots,X_s$ be fixed irreducible positive-dimensional
projective varieties over $\Qbar$, with $s\ge2$, and fix embeddings
$X_\ell\subseteq\mathbb P^{N_\ell}$ given by very ample bundles
$L_\ell$.  For positive integral weights $a_\ell$, put
$N_a=\boxtimes_\ell L_\ell^{a_\ell}$.
Let $\mathcal M$ be a nef bundle on $\prod_\ell X_\ell$ and
$\mathcal P=N_a^{t_1}$, with its full monomial coordinate system $\Xi$.
Suppose $\mathcal P\otimes\mathcal M^{-1}$ is generated by $J$ sections
$\mathcal Z$, whose images under an injection
\[
 \mathcal P\otimes\mathcal M^{-1}\hookrightarrow N_a^{t_2}
\]
are polynomials of multidegree $t_2a$ with joint coefficient height at
most $\sum_\ell a_\ell\delta_\ell$.  Assume the injection is an
isomorphism on a nonempty open set $U^0$.  Here $t_1,t_2,J$ are positive
integers and $\delta_\ell\ge1$.

Let $x=(x_\ell)\in U^0(\Qbar)$.  Suppose that, for every irreducible
subproduct $Z=\prod_\ell Z_\ell$ with $x_\ell\in Z_\ell\subseteq X_\ell$,
\begin{equation}\label{eq:dill-intersections}
 (\mathcal M^{\dim Z}\cdot Z)
 \ge \theta^{-1}\prod_\ell
       (\deg Z_\ell)^{-\omega}a_\ell^{\dim Z_\ell},
\end{equation}
where $\theta\ge1$ and $\omega\ge-1$ are integers.  Put
\[
 \begin{split}
 u_0&=\sum_\ell\dim X_\ell,\\
 \Lambda&=\theta(2t_1u_0)^{u_0}
          \max_\ell(N_\ell+1)\prod_\ell\deg X_\ell,\\
 \psi(u)&=\prod_{j=u+1}^{u_0}((3+\omega)j+1),
 \qquad c_1=c_2=\Lambda^{\psi(0)},\\
 c_3^{(\ell)}&=\Lambda^{2\psi(0)}(Jt_2)^{u_0}
                      (h(X_\ell)+\delta_\ell).
 \end{split}
\]
If $a_\ell\ge c_2a_{\ell+1}$ for $\ell<s$ and
$h(x_\ell)\ge c_3^{(\ell)}$ for every $\ell$, then
\begin{equation}\label{eq:dill-height}
 \sum_\ell a_\ell h(x_\ell)
 \le c_1\bigl(h(\Xi(x))-h(\mathcal Z(x))\bigr).
\end{equation}
\end{theorem}

This is \cite[Theorem~1.1]{Dill2020} with the product as its own proper
model and with the identity injection
$\mathcal P\hookrightarrow N_a^{t_1}$.  In that theorem the intersection
is taken on the closure of the inverse image of $Z\cap U^0$; here this
closure is $Z$, because $Z$ is irreducible and contains $x\in U^0$.
The constants do not involve the complement of $U^0$.  In particular,
the statement can be applied to different such open sets with the same
numerical parameters.

The crucial distinction is that $c_1$ depends only on the geometric data,
$t_1$, $\theta$, and $\omega$.  The number $J$ of generators, their
coefficient heights, and $t_2$ enter only the height thresholds
$c_3^{(\ell)}$.  We will fix $c_1$ before choosing a rational
approximation; the height thresholds may be fixed afterward.

\subsection{The all-high contradiction}

\begin{theorem}\label{thm:all-high}
Theorem~\ref{thm:sequence} holds for all-high sequence data.
\end{theorem}

\begin{proof}
We first fix a positivity neighborhood and the constant in Dill's
inequality. For an arbitrary approximation accuracy, we then balance the
height scales by integer multipliers and construct the required bundles.
The final comparison will determine how small the accuracy must be.

Suppose that all-high data violate Theorem~\ref{thm:sequence}, and
choose a counterexample with $\dim\cA$ minimal.
The projection exclusion and Proposition~\ref{prop:tangent} then apply.
Every parent $Y_\ell$ has positive dimension: if it were a point,
the positive normalized lower bound \eqref{eq:sequence-nondegenerate},
applied to the identity homomorphism, would contradict $H_{i\ell}\to\infty$.
Choose a closed rational box $\mathcal B$ which contains $\Phi$ in its
interior and lies in the neighborhood supplied by
Proposition~\ref{prop:tangent}.  After discarding finitely many indices,
$P_Z(g)>0$ on this box for every subproduct through $y_i$ in the fixed
parents.  Lemma~\ref{lemma:intersection-lower-bound} gives constants
$\theta,\omega_1$ valid for all these products.

If there is only one factor, use two subsequences $n_i,m_i\to\infty$
with $H_{n_i}/H_{m_i}\to0$.  They exist by choosing
$H_{m_i}\ge i^2H_{n_i}$.  On the two copies use the maps $g\oplus g$
and the external product target polarization.  For subvarieties
$Z_1,Z_2$ through the corresponding points,
\[
 P_{Z_1\times Z_2}(g\oplus g)
   =\binom{\dim Z_1+\dim Z_2}{\dim Z_1}P_{Z_1}(g)P_{Z_2}(g).
\]
Thus the required bound holds with $\theta^2$ and the same $\omega_1$.
The block diagonal limiting projector still sends the normalized
point pairs to zero.  From now on there are at least two factors.
In the duplicated case all approximations remain in this block diagonal
family; no new minimality argument in the larger ambient is needed.
We relabel the doubled data and replace $\theta$ by $\theta^2$ in this
case, retaining the notation already introduced.

Choose a positive integer $t_1$, fixed independently of the rational
approximation, so large that
\[
 t_1\sum_\ell\operatorname{pr}_\ell^*\omega_{L_\ell}-g^*\omega_L
\]
is positive definite for every $g\in\mathcal B$ (or for the indicated
block diagonal family).  Compactness of the box makes this possible.
With the parent embeddings and $\theta,\omega_1,t_1$ fixed,
Theorem~\ref{thm:dill-form} has a fixed multiplicative constant $c_1$.

Let $0<\eta<1/2$.  The following construction is available for every
such $\eta$, with $c_1$ already fixed.  Choose a rational homomorphism
$f/e$ in the box with $\|f/e-\Phi\|\le\eta$, where $e>0$ is even
and $f$ is an integral homomorphism divisible by two.

We want the multiplied points in every factor to have the same height
scale: $v_{i\ell}\sqrt{H_{i\ell}}$ should be close to one number $R_i$.
At the same time, the multiplication maps must have polynomial formulas
whose coefficient heights grow at most quadratically in $v_{i\ell}$.
Both requirements are met by a dyadic family.  Choose an integer
$N>2/\eta$ and put $U=\{N,N+1,\ldots,2N\}$.  The numbers $u2^k$,
with $u\in U$ and $k\ge0$, approximate every real number at least $N$
to relative error at most $\eta$.  Taking
$R_i\ge N\max_\ell\sqrt{H_{i\ell}}$, choose
\begin{equation}\label{eq:dyadic-multipliers}
 v_{i\ell}=u2^k,\qquad u\in U,\quad k\ge0,
\end{equation}
so that
\[
 1-\eta\le\frac{v_{i\ell}\sqrt{H_{i\ell}}}{R_i}\le1+\eta.
\]
We verify the coefficient-height requirement below.

On the ambient product set
\[
 \mathcal Q_0=\left(\boxtimes_\ell L_\ell^{t_1e^2}\right)
                           \otimes(f^*L)^{-1}.
\]
The choice of $t_1$ makes its invariant curvature positive definite,
hence makes this bundle ample
\cite[Theorem~1.1.4]{GenestierNgo2009}.  Writing $e=2e_0$ and
$f=2f_0$, symmetry identifies it with the fourth power of the ample
bundle
\[
 \left(\boxtimes_\ell L_\ell^{t_1e_0^2}\right)
                           \otimes(f_0^*L)^{-1}.
\]
It is therefore globally generated
\cite[Theorem~2.2.3]{GenestierNgo2009}. Fix a generating system
$z_1,\ldots,z_J$ for $\mathcal Q_0$.

Let $[v_i]$ denote factorwise multiplication by the chosen integers, and
define
\[
 a_{i\ell}=e^2v_{i\ell}^2,\qquad
 N_{a_i}=\boxtimes_\ell L_\ell^{a_{i\ell}},\qquad
 \mathcal M_i=(f\circ[v_i])^*L\big|_{\prod Y_\ell},\qquad
 \mathcal P_i=N_{a_i}^{t_1}.
\]
The bundle $\mathcal M_i$ is nef.  Symmetry identifies
$\mathcal P_i\otimes\mathcal M_i^{-1}$ with
$[v_i]^*\mathcal Q_0|_{\prod Y_\ell}$, which is generated by the
pullbacks of the $z_j$.

Choose a coordinate section $\sigma$ of the target $L$ which is
nonzero at $f([v_i]y_i)$.  Multiplication by
$(f\circ[v_i])^*\sigma$ gives the injection
\[
 \mathcal P_i\otimes\mathcal M_i^{-1}
       \hookrightarrow N_{a_i}^{t_1}.
\]
It is an isomorphism on the nonvanishing open set of that section,
which contains $y_i$.  Thus we may take $t_2=t_1$ in
Theorem~\ref{thm:dill-form}.  Before pullback, each product
$z_jf^*\sigma$ is a section of
$\boxtimes_\ell L_\ell^{t_1e^2}$.
Projective normality of the ambient factor embeddings, together with
the external-product formula for global sections, represents these
sections by multihomogeneous polynomials of multidegree
$(t_1e^2,\ldots,t_1e^2)$.  There are only finitely many systems, one
for each choice of target coordinate $\sigma$.

We next check that the dyadic choice \eqref{eq:dyadic-multipliers}
gives the uniform coefficient bound required by Dill.
Symmetry and projective normality give polynomial coordinate systems
of degree four for $[2]$ and of degree $u^2$ for each $[u]$, $u\in U$.
Their compositions give the required degree $v_{i\ell}^2$ formulas.

Fix a number field containing the ambient data, $f$, and the chosen
generating systems.  At a finite
place use the maximum coefficient norm, and at an infinite place use
the sum of absolute coefficient values.  If $b_w(k)$ is the logarithm
of the maximum of one and the coefficient norm after $k$ iterations
of the degree-four system, submultiplicativity gives
\[
 b_w(k+1)\le4b_w(k)+C_w.
\]
Thus $b_w(k)=O_w(4^k)$.  Composition with one of the finitely many
$[u]$ gives $O_w((u2^k)^2)$.  At finite places the constants vanish
outside a fixed finite set.  Substitution in the finitely many fixed
systems representing $z_jf^*\sigma$ now bounds their joint
coefficient height by
\[
 C_{f,e,U}\sum_\ell a_{i\ell}.
\]
Consequently we may fix $\delta_\ell=\max(1,C_{f,e,U})$ for every
$\ell$, independently of $i$ and of the chosen nonvanishing coordinate.
The generator count $J$ is also fixed.  Dependence on the approximation
and on $U$ changes only the height thresholds.

For every subproduct $Z$ through $y_i$, the scaling identity for $P_Z$
and \eqref{eq:intersection-lower-bound} give
\[
 (\mathcal M_i^{\dim Z}\cdot Z)
  =P_Z(f/e)\prod_\ell a_{i\ell}^{\dim Z_\ell}
  \ge\theta^{-1}\prod_\ell
       (\deg Z_\ell)^{-\omega_1}a_{i\ell}^{\dim Z_\ell}.
\]
This verifies the last geometric hypothesis of Dill's theorem.

\medskip\noindent
\emph{Height comparison and choice of accuracy.}

All geometric and coefficient conditions have now been checked for
each fixed $\eta$.  The balancing of the multipliers will make the
source height large and the image height small after division by
$e^2R_i^2$.

Let $\Xi_i$ be the full monomial coordinates for $\mathcal P_i$, and
let $\mathcal Z_i$ be the pulled-back generating system just constructed.
Then
\begin{equation}\label{eq:high-height-comparison}
 h(\Xi_i(y_i))-h(\mathcal Z_i(y_i))
  =\widehat h_L(f([v_i]y_i))
       +O_{f,e}\left(1+\sum_\ell a_{i\ell}\right).
\end{equation}
Indeed, the monomial height is exactly
$t_1\sum_\ell a_{i\ell}h(y_{i\ell})$.
Multiplying all generator values by the same nonzero injection section
does not change their projective height.  Their height is therefore
the fixed generator height of $\mathcal Q_0$ at $[v_i]y_i$.
As $\mathcal Q_0$ is symmetric, this equals
\[
 t_1\sum_\ell a_{i\ell}\widehat h_{L_\ell}(y_{i\ell})
      -\widehat h_L(f([v_i]y_i))+O_{f,e}(1).
\]
The bounded difference between projective and canonical heights on
each fixed factor proves \eqref{eq:high-height-comparison}.

By the high-scale condition applied to the identity, there are constants
$\kappa_\ell>0$, independent of the approximation, such that
\[
 h(y_{i\ell})\ge\kappa_\ell H_{i\ell}
\]
for all sufficiently large $i$.  The normalized vectors $q_i$ are
bounded.  Operator-norm comparison for real homomorphisms therefore
gives a constant $C$ independent of $\eta$, $f/e$, and the multipliers:
the maps $f/e$ lie in the fixed box, and the balancing ratios lie in
$[1/2,3/2]$.
Let $D_i$ be the diagonal real endomorphism with entries
$v_{i\ell}\sqrt{H_{i\ell}}/R_i$.  In the real height space,
\[
 \frac{f([v_i]y_i)}{eR_i}=(f/e)D_iq_i.
\]
Writing this as the sum of $\Phi q_i$,
$(f/e-\Phi)D_iq_i$, and $\Phi(D_i-1)q_i$ gives
\begin{equation}\label{eq:high-small-image}
 \frac{\sqrt{\widehat h_L(f([v_i]y_i))}}{eR_i}
   \le C\eta+o(1).
\end{equation}

For the weights, separation of successive scales gives
\[
 \frac{a_{i\ell}}{a_{i,\ell+1}}\longrightarrow\infty,
 \qquad
 \frac{\sum_\ell a_{i\ell}}{e^2R_i^2}
       \le (1+\eta)^2\sum_\ell H_{i\ell}^{-1}\longrightarrow0.
\]
Moreover
\begin{equation}\label{eq:high-large-source}
 \frac{\sum_\ell a_{i\ell}h(y_{i\ell})}{e^2R_i^2}
       \ge(1-\eta)^2\sum_\ell\kappa_\ell
       \ge\frac14\sum_\ell\kappa_\ell.
\end{equation}
Once $\eta$, the approximation, and $U$ are fixed, all the height
thresholds in Theorem~\ref{thm:dill-form} are fixed.  The projective
heights tend to infinity and the weight ratios cross the fixed ratio
threshold, so \eqref{eq:dill-height} applies for all sufficiently
large $i$.  Dividing it by $e^2R_i^2$ and using
\eqref{eq:high-height-comparison}--\eqref{eq:high-large-source} yields
\[
 \frac14\sum_\ell\kappa_\ell
       \le c_1 C^2\eta^2+o(1).
\]
The construction and this inequality hold for every $0<\eta<1/2$,
whereas $c_1$, $C$, and the $\kappa_\ell$ are independent of $\eta$.
Choose $\eta$ so that $c_1C^2\eta^2<\frac14\sum_\ell\kappa_\ell$.
Keeping its associated $f,e,U$ fixed and letting $i$ tend to infinity
gives the contradiction.  All constants depending on this approximation
enter only fixed thresholds or errors which vanish after normalization.
\end{proof}
